%% 03-results.tex — Three bullet results

\section{What the Evidence Shows}

Applied to all fifty states across the 2000, 2010, and 2020 Census cycles,
the algorithm produces the following results:

\bigskip

\noindent
\begin{tabular}{@{}>{\bfseries\large}p{0.08\textwidth} p{0.86\textwidth}@{}}
+22\% &
More compact than enacted 2020 maps, measured by the Polsby-Popper
ratio of district area to perimeter squared, averaged across all 435 districts.
Compact districts follow rivers and county lines rather than
twisting through neighborhoods to sort voters.
\textit{(Paper B.2~\cite{dellaLibera2026edgeWeighted})}\\[10pt]

7D / 7R &
North Carolina---the state that has generated more federal gerrymandering
litigation than any other in the past decade---receives a
seven-Democratic, seven-Republican delegation under the algorithm,
closely matching the state's evenly divided electorate.
No partisan data was used.
\textit{(Paper B.11~\cite{dellaLibera2026ncRegions})}\\[10pt]

15--20 min &
Any citizen with a laptop and an internet connection can download
the open-source \texttt{bisect} tool, fetch public Census data,
run the Vermont verification fixture, and confirm the result themselves
in 15--20 minutes end-to-end including data download.
Full fifty-state replication takes approximately 18 minutes for one
census year, or two to four hours for all three census years.
\textit{(Replication Package A.4)}\\
\end{tabular}

\bigskip

\bigskip

\noindent
\begin{tabular}{@{}>{\bfseries\large}p{0.08\textwidth} p{0.86\textwidth}@{}}
50 states &
The algorithm scales from Vermont's one congressional district to New Hampshire's
400-member state House.
All 50 state lower chambers have been redistricted; 39 require block-group
resolution (finer than the census-tract standard used for Congress).
State house maps are 7.5\% more compact than congressional maps by construction---
smaller districts follow geographic features more precisely.
\textit{(Papers F.1, F.3)}\\[10pt]

NestSection &
For states whose senate-to-house seat ratio has a common factor (40 of 49
bicameral states), the NestSection algorithm produces senate districts
that are exact unions of house districts---ensuring that every senate
constituent lives in exactly one house district, eliminating the
double-gerrymandering problem that plagues bicameral states today.
\textit{(Paper F.2)}\\
\end{tabular}

\bigskip

\bigskip

\noindent
\begin{tabular}{@{}>{\bfseries\large}p{0.08\textwidth} p{0.86\textwidth}@{}}
0.21 &
The \emph{Representation Quality Index} (Track O) synthesises four
democratic outcome dimensions into a single score: algorithmic plans score
$0.62$ vs.\ $0.41$ for enacted gerrymanders.
The specific advantages: $+23$pp more competitive districts, $+1.1$pp
higher midterm voter turnout, $-0.07$ less extreme DW-NOMINATE scores,
and $14$ fewer minutes mean constituent-to-representative travel time.
\textit{(Papers O.1--O.5)}\\[10pt]

\$50K &
Total implementation cost for a state redistricting office to run the
full \texttt{bisect} pipeline: census data download (\$0, public domain),
compute (18 minutes, $<\$1$ on cloud), staff review time ($\approx\$50{,}000$).
This compares to Iowa's \$3~million nonpartisan process and California's
\$35~million commission process.
Algorithmic redistricting is not only better --- it is cheaper by $60$--$700\times$.
\textit{(Paper P.5)}\\
\end{tabular}

\bigskip

\noindent Additional findings:
the algorithm produces \textbf{more majority-minority districts}
than enacted plans in states above the 42\% minority-population threshold
identified in Paper D.1~\cite{dellaLibera2026vraCompliance},
and reduces the partisan efficiency gap by \textbf{62\%} compared to
enacted plans~\cite{dellaLibera2026efficiencyGap}---without any
partisan fairness criterion in the algorithm specification.\footnote{%
The algorithm produces near-zero efficiency gaps as a byproduct of geometric
neutrality (mean $|EG|=0.04$ vs.\ $0.08$ for enacted maps).
In some geographically sorted states a small systematic gap persists due to
urban concentration (Paper C.5).}
